\section{Fractional densities and a harmonic infinity bound}\label{tra:fractional-density}

This section gives an alternative proof of the balanced compatibility
bound using fractional integrability and a softened survival potential.
Theorem~\ref{tra:duality:weighted} already implies the compatibility
estimate below for sufficiently large sizes. The proof here obtains its
coefficient integrability by spectral bins, without inverse
Hausdorff--Young, and carries that bound through the entropy exposure.
Its separate sparse result is an infinity-norm estimate on harmonic
tuples, with an explicit constant and the full range $k\le n/(2d)$.

One sweep is $d$ physical shuffles. All regular traces are unnormalized.
We use $T_n$ and $Q_n=T_n^*T_n$ as in Section~\ref{sec:prelim}.

\label{tra:fractional-density:density-normalization}
We will use the left regular representation of \(S_n\). For a finite group \(G\) in general, write \(\mathcal L(g)\) for the left regular operators. We identify group algebra elements with operators in this representation, calling \(f\) the density of \(Z\) if
\[
 Z=\frac{1}{|G|}\sum_{g\in G} f(g)\mathcal L(g)
\]
(where \(f\) need not be a probability density). With ordinary, \emph{unnormalized} trace, we have
\begin{equation}
 f(g)=\operatorname{Tr}\mathcal L(g^{-1})Z,\qquad
 \mathbb E_{g\ \mathrm{unif.}} |f(g)|^2=\operatorname{Tr}(Z^*Z).
 \label{tra:fractional-density:eq1}
\end{equation}
For \(S_n\), let \(\Pi_i\) be the average operator for the random swap layer at bit \(i\). This is an orthogonal projection, since it averages over the subgroup generated by all the disjoint transpositions in the layer. Set
\[
 T_n=\Pi_d\cdots\Pi_1,\qquad Q_n=T_n^* T_n,\qquad
 M_n(p)=\operatorname{Tr}(Q_n^p)\quad (p>0).
\]
We have \(0\le Q_n\le I\), and \(M_n(p)\ge 1\) by the constants. We use the notation \(T_n,Q_n,M_n\) also with other powers of two as the size (using that size's own scan).

\label{tra:fractional-density:power-log-majorization-and-mixing}
For later comparison, a positive integer \(L\) satisfies
\begin{equation}
 n!\sum_g \Pr(g\ \text{is the movement in }L\ \text{scans})^2
  =\operatorname{Tr}\big((T_n^L)^*T_n^L\big)\ \le\ M_n(L).
 \label{tra:fractional-density:eq3}
\end{equation}
The inequality is Schatten H\"older, as used in
Proposition~\ref{e:nonnormal-moment-conversion}. A positive-square moment of
order $p$ is a singular moment of order $2p$, so
$M_n(p)\le1+1/8$ supplies Proposition~\ref{e:nonnormal-moment-conversion}
with singular exponent $2p$ and remainder $1/8$.

The positive-power comparison, Lemma~\ref{lem:positive-power-comparison},
with $P=4r$ and $q=4$, also gives the form used below:
\begin{equation}
 \operatorname{Tr}(B^{1/2}AB^{1/2})^{2r}
 \le\operatorname{Tr}(A^rB^rA^rB^r)
 \qquad(A,B\ge0,\ r\ge1).
 \label{tra:fractional-density:eq4}
\end{equation}

\subsection{The induction step and rough moment}

\label{tra:fractional-density:induction-budgets}
All logarithms in the proof below are natural logs. Fix the constants
\[
 \alpha=\frac{1}{64},\qquad \beta=\frac{1}{256},\qquad c_0=\frac18,
\]
and let
\[
 a=2^{\lceil d/2\rceil},\qquad b=2^{\lfloor d/2\rfloor},\qquad
 n=ab,\qquad \delta=d^{-1/2}.
\]

\begin{proposition}\label{tra:fractional-density:step}

\label{tra:fractional-density:step-and-rough-moment}
For all sufficiently large \(d\), if a real \(q\ge16/\beta\) satisfies
\begin{equation}
 M_b(q), M_a(q) \ \le\ 1+c_0,
 \label{tra:fractional-density:eq5}
\end{equation}
then \(p_0=(1+\delta)q\) satisfies
\begin{equation}
 M_n((1+\delta)p_0)\le 1+c_0
 \label{tra:fractional-density:eq6}
\end{equation}
and
\begin{equation}
 M_n(p_0)\le \exp(n^{1-\alpha}).
 \label{tra:fractional-density:eq7}
\end{equation}
The lower threshold for $d$ is absolute and independent of $q$.

\end{proposition}

We first prove the rough estimate \eqref{tra:fractional-density:eq7}
by splitting the array. Its sublinear exponent will control the
large-dimensional representations in \eqref{tra:fractional-density:eq6};
small levels will be bounded directly on tuples of positions.

\label{tra:fractional-density:balanced-operator-split}
\textbf{Rough bound via a balanced split.} View positions as an \([a]\)-by-\([b]\) array with row index made from the last \(\lceil d/2\rceil\) bits and column index from the first \(\lfloor d/2\rfloor\) bits (using \([m]=\{1,\ldots,m\}\)). The scan first operates by horizontal permutations in
\[
 H=(S_b)^a
\]
and then by vertical permutations in
\[
 V=(S_a)^b
\]
where these are the subgroups preserving each row and each column, respectively. In fact the first part of the scan just does independent size-\(b\) scans, one in each row, and then the second part does independent size-\(a\) scans, one in each column. Write the average operators of those two parts as \(K_H,K_V\), so \(T_n=K_V K_H\). Set
\[
 A=K_H K_H^*,\qquad B=K_V^* K_V,\qquad r=p_0/2.
\]
Since \(Q_n=K_H^* B K_H\), it has the same eigenvalues as \(B^{1/2} A B^{1/2}\). Thus \eqref{tra:fractional-density:eq4} gives
\begin{equation}
 M_n(p_0)\le \operatorname{Tr}(A^r B^r A^r B^r).
 \label{tra:fractional-density:eq8}
\end{equation}

\label{tra:fractional-density:coefficient-probability-reduction}
The line factors of $A^r$ and $B^r$ are
$(T_bT_b^*)^r$ and $(T_a^*T_a)^r$, respectively. Let their coefficient
densities be $f_R,f_D$. Their squared masses are $M_b(p_0),M_a(p_0)$,
both in $[1,2]$ by \eqref{tra:fractional-density:eq5} and $p_0\ge q$.

For $h=(h_i)\in H$ and $v=(v_j)\in V$, the entropy exposure uses arrays
\[
 X_{ij}=h_i(j)\in[b],\qquad Y_{ij}=v_j^{-1}(i)\in[a],
 \qquad
 \mathcal E=\{(X_{ij},Y_{ij})\text{ are all distinct}\}.
\]
This is the domain $hv\in VH$ in the proof of
Lemma~\ref{lem:coefficient-compatibility}: at the intermediate cell
$(i,j)$, after $v$ and before $h$, the input row is $Y_{ij}$ and
the output column is $X_{ij}$. Opposite-order routing is possible
exactly when each input row sends one card to each output column.
Define the product probability law $\nu$ on these arrays by the
independent line densities
\[
 \rho_R(x)=\frac{|f_R(x)|^2}{M_b(p_0)},\qquad
 \rho_D(y)=\frac{|f_D(y^{-1})|^2}{M_a(p_0)}.
\]
The same lemma, using the equivalent $hv\in VH$ form of its
four-factor calculation, gives
\begin{equation}
 M_n(p_0)\le\frac{n!}{|H||V|}
      M_b(p_0)^a M_a(p_0)^b\,\nu(\mathcal E).
 \label{tra:fractional-density:eq10}
\end{equation}
We estimate the last probability, keeping the estimate robust under the density changes in \(\nu\).

\label{tra:fractional-density:fractional-spectral-density}
First consider a single row or column law \(\nu_*\) in this product, on \(S_m\) with \(m=b\) or \(a\). For the row case let \(T=(T_m T_m^*)^q\); for the column case let \(T=(T_m^* T_m)^q\). By hypothesis \(0\le T\le I\) and \(\operatorname{Tr} T\le 2\) in the regular representation for this \(m\). We deal with the density of \(F=T^{(1+\delta)/2}\), evaluated on permutations for the row case and on their inverses for the column case. Divide the positive eigenvalues of \(T\) into intervals
\[
 (e^{-l-1},e^{-l}]\qquad (l=0,1,\ldots),
\]
and correspondingly write \(F=\sum_l F_l\) using spectral parts (only finitely many nonzero). Then
\[
 \operatorname{Tr} F_l\le 2 e^{(1-\delta)(l+1)/2},
 \qquad
 \operatorname{Tr} F_l^2\le 2 e^{-\delta l}.
\]
Write \(f_l\) for the associated density (evaluated on inverses in the column case). By \eqref{tra:fractional-density:eq1},
\[
 \|f_l\|_\infty\le \operatorname{Tr} F_l\le 2e^{(l+1)/2},\qquad
 \|f_l\|_2^2\le 2e^{-\delta l},
\]
where the first inequality follows because \(F_l\) is positive semidefinite and the \(\mathcal L(g)\) are unitary. Norms for functions here are with uniform measure. Thus
\[
 \|f_l\|_{2+\delta}^{2+\delta}
 \le C e^{-\delta l/2}
\]
for an absolute \(C\), taking \(d\) large so \(0<\delta\le 1\). Adding the norms gives
\[
 \left\|\sum_l f_l\right\|_{2+\delta}\le C'/\delta
\]
by summing a geometric series, with \(C'\) absolute. Therefore the density \(\rho\) of \(\nu_*\), which equals the square modulus of \(\sum_l f_l\) divided by a number at least 1, satisfies
\begin{equation}
 \mathbb E_{U_m}\rho^{1+\delta/2}\le (C'/\delta)^{2+\delta}
 \label{tra:fractional-density:eq11}
\end{equation}
where \(U_m\) denotes uniform on \(S_m\).

\label{tra:fractional-density:density-relative-entropy-comparison}
We use relative entropy \(D(P\|R)=\mathbb E_P\log(P/R)\) for probability laws, infinite if \(P\) is not supported on the support of \(R\). From \eqref{tra:fractional-density:eq11}, with
\[
 \epsilon=\frac{\delta}{2+\delta},
\]
any law \(\mu\) on \(S_m\) satisfies
\begin{equation}
 D(\mu\|\nu_*)\ \ge\ \epsilon D(\mu\|U_m)-O(\log d).
 \label{tra:fractional-density:eq12}
\end{equation}
In fact we use the elementary variational inequality
\begin{equation}
 D(P\|R)\ \ge\ \mathbb E_P W-\log \mathbb E_R e^W
 \label{tra:fractional-density:eq13}
\end{equation}
for potentials \(W\) (also allowing \(-\infty\) on points outside the support of \(P\) with the exponential expectation positive). This follows by changing the reference measure by the exponential tilt and using nonnegativity of relative entropy. In \eqref{tra:fractional-density:eq12}, assuming finite divergence, apply \eqref{tra:fractional-density:eq13} with \(P=\mu\), \(R=U_m\), and \(W=(1+\delta/2)\log\rho\), and use \(D(\mu\|\nu_*)=D(\mu\|U_m)-\mathbb E_\mu\log\rho\). This gives \eqref{tra:fractional-density:eq12} because of \eqref{tra:fractional-density:eq11} and \(\delta=d^{-1/2}\).

\begin{lemma}\label{tra:fractional-density:compatibility}

\label{tra:fractional-density:compatibility-probability}
The compatibility probability satisfies
\begin{equation}
 \nu(\mathcal E)
 \ \le\ \exp\left(-n+O\left(n a^{-1/16}\log a+(a+b)\log d\right)\right).
 \label{tra:fractional-density:eq14}
\end{equation}

\end{lemma}\begin{proof}\label{tra:fractional-density:compatibility-proof}

\label{tra:fractional-density:entropy-decomposition-and-reveal}
Consider any law \(\zeta\) on arrays supported on \(\mathcal E\), assuming \(D(\zeta\|\nu)<\infty\). Write \(X_i=X_{i,\cdot}\) for the row permutations and \(Y_j=Y_{\cdot,j}\) for the column permutations (using single subscripts for these permutations). We use \(\zeta_X,\zeta_{X_i},\zeta_{Y_j}\), etc., for the laws of parts of the arrays under \(\zeta\). Put
\[
 D_X=D(\zeta_X\|U_b^{\otimes a}),\qquad
 D_Y=\sum_{j\in[b]}D(\zeta_{Y_j}\|U_a),\qquad
 J=\mathbb E_{X\sim\zeta_X}D\Big(\zeta_{Y\mid X}\ \Big\|\ \bigotimes_{j\in[b]}\zeta_{Y_j}\Big).
\]
By \eqref{tra:fractional-density:eq12},
\begin{equation}
 D(\zeta\|\nu)
 \ \ge\ J+\epsilon(D_X+D_Y)-O((a+b)\log d).
 \label{tra:fractional-density:eq15}
\end{equation}
Indeed, decomposing the joint divergence against the product law \(\nu\), we get \(J\), plus \(D(\zeta_X\| \bigotimes_i \nu_{X_i})\), plus the divergences of the \(\zeta_{Y_j}\) against their individual laws in \(\nu\). For the \(X\) term, decompose further into \(D(\zeta_X\| \bigotimes_i\zeta_{X_i})\) and the sum of individual marginal divergences. Applying \eqref{tra:fractional-density:eq12} to the individual divergences and keeping at least an \(\epsilon\) fraction of \(D(\zeta_X\| \bigotimes_i\zeta_{X_i})\ge 0\) gives \eqref{tra:fractional-density:eq15}.

To lower bound \(J\), consider drawing the arrays from \(\zeta\), with \(X\) revealed first. Take independent uniform priorities \(t_j\in[0,1]\), independent also of the arrays, and expose \(Y\) column by column, in order of \emph{decreasing} \(t_j\). Within each column reveal its entries in increasing order of \(i\). For any fixed priorities (ignoring null events of ties), the chain formula for relative entropy gives \(J\) as a sum of expected conditional divergences, one per cell \((i,j)\). In that cell the divergence is between:
\begin{itemize}
\item the conditional distribution of \(Y_{ij}\) under \(\zeta\), given \(X\), earlier columns, and earlier entries in column \(j\);
\item the law \(u\) with
\end{itemize}
\[
 u_y=\zeta_{Y_j}(Y_{ij}=y\mid Y_{1j},\ldots,Y_{i-1,j}),\qquad y\in[a].
\]
The second (reference) law depends just on previous entries in the \emph{same} column. Though we suppress its cell and history indices on \(u,u_y\), we always take conditionals at histories arising under \(\zeta\). Since we know \(X\), the first distribution is supported on candidates that are not \emph{blocked}, where \(y\) is blocked if \((X_{ij},y)\) has already occurred as a pair in an earlier column.

Put
\[
 g(j,x)=|\{i': X_{i'j}=x\}|.
\]
Call the cell bad if \(g(j,X_{ij})>a^{1/8}\), good otherwise. Call a candidate \(y\) high if \(u_y>a^{-1/4}\), low otherwise. Let \(R_X\) count bad cells, and \(R_Y\) cells whose actual \(Y_{ij}\) is high. These counts depend on the outcome arrays (and law \(\zeta\) for defining high), but not the priorities.

\label{tra:fractional-density:low-candidate-potential}
For a good cell in the chain formula, apply \eqref{tra:fractional-density:eq13} to the conditional divergence, using \(\eta=a^{-1/16}\) and this potential:
\[
 W(y)=
 \begin{cases}
 0, & y\ \text{high},\\
 -\infty, & y\ \text{low and blocked},\\
 -\log(t_j+\eta), & y\ \text{low and not blocked}.
 \end{cases}
\]
On bad cells just take \(W(y)=0\). These choices depend only on what is given before the cell, besides the priorities and the law. The actual choice under \(\zeta\) is not blocked and has \(u_y>0\) almost surely, so the reference normalization in \eqref{tra:fractional-density:eq13} is positive as needed.

For the log normalization term of \eqref{tra:fractional-density:eq13}, condition now on the entire outcome arrays \((X,Y)\) and on \(t_j\), and average over the other priorities. Assume the cell is good and write \(x=X_{ij}\). For each \(y\), including each low candidate, the pair \((x,y)\) occurs exactly once in the outcome arrays, since \(\mathcal E\) holds and there are \(ab\) cells. If its occurrence is in another column, the probability the candidate is not blocked is \(t_j\), since the priority of that column is uniform independent. The candidates occurring in the current column with \(x\) number at most \(g(j,x)\), so the total \(u\) mass of those among them which are low is at most \(a^{1/8} a^{-1/4}\). Writing \(Z=\sum_y u_y e^{W(y)}\), this gives
\[
 \mathbb E\big[ Z\mid X,Y,t_j\big]
 \ \le\ \sum_{y\ \mathrm{high}} u_y
       +\frac{t_j}{t_j+\eta}\sum_{y\ \mathrm{low}} u_y
       +\frac{a^{1/8}a^{-1/4}}{\eta}
 \ \le\ 1+a^{-1/16}.
\]
This computation uses that \(u\) is fixed in this conditioning, as the within-column order is fixed. On a bad cell \(Z=1\). We get in particular \(\mathbb E\log Z\le a^{-1/16}\) by Jensen's inequality.

For the expected-potential term in \eqref{tra:fractional-density:eq13}, averaging over histories under \(\zeta\) amounts to using the actual score \(W(Y_{ij})\). For each good cell whose actual choice is low, since good/low status for the actual choice is independent of the priorities, its score with the priorities averaged is
\[
 \int_0^1 -\log(t+\eta)\,dt
 =1-(1+\eta)\log(1+\eta)+\eta\log\eta
 =1-O(a^{-1/16}\log a),
\]
and it is 0 on the other cells. The displayed integral lies between 0 and 1 for sufficiently large \(a\). Since the chain formula for \(J\) holds for each fixed priority vector with distinct entries, averaging our lower bounds over priorities and summing all cells gives
\begin{equation}
 J\ \ge\ n-\mathbb E_\zeta[R_X+R_Y]-O(n a^{-1/16}\log a).
 \label{tra:fractional-density:eq16}
\end{equation}

\label{tra:fractional-density:high-atom-and-clump-charge}
The counts lost in \eqref{tra:fractional-density:eq16} can be paid for by the deficits \(D_X,D_Y\). First, for \(D_Y\), use its chain formula for the marginal column laws against uniform permutations, within each column in increasing \(i\). While the number remaining to be revealed in the column (including the current cell) is at least \(\sqrt a\), the high subset has uniform conditional probability at most \(a^{-1/4}\), since it has size at most \(a^{1/4}\). Use \eqref{tra:fractional-density:eq13} with potential \((\log a)/8\) on the high subset and 0 outside, to see that the conditional divergence at this cell, from \(u\) against the uniform choice among those remaining, is at least
\[
 \frac{\log a}{8}\sum_{y\ \mathrm{high}} u_y - a^{-1/8}.
\]
Skipping fewer than \(\sqrt a\) cells at the end of each column (their conditional divergences are nonnegative), taking marginal expectations, and summing gives
\begin{equation}
 D_Y\ \ge\ \frac{\log a}{8}\big(\mathbb E_\zeta R_Y - n a^{-1/2}\big)-n a^{-1/8}.
 \label{tra:fractional-density:eq17}
\end{equation}
For \(D_X\), we claim for large \(a\),
\begin{equation}
 D_X\ \ge\ \frac{\log a}{16}\mathbb E_\zeta R_X-1.
 \label{tra:fractional-density:eq18}
\end{equation}
To verify this let \(\gamma=(\log a)/16\) and \(h_0=a^{1/8}\). Under \(U_b^{\otimes a}\), the reference law for \(D_X\), the probability of any specification of \(z\) different cells to prescribed values in \(X\) is at most \((e/b)^z\). For a compatible specification this follows row by row because for \(z_i\) cells in one row its probability is \(1/(b(b-1)\cdots(b-z_i+1))\), at most \((e/b)^{z_i}\) by \(b!\ge (b/e)^b\) (the geometric mean in the descending product for the row is at least that from \(b!\)). Incompatible specifications have probability zero. Also,
\[
 e^{\gamma R_X}
  =\prod_{j,x\in[b]}
       \left(1+\big(e^{\gamma g(j,x)}-1\big)\mathbf 1_{\{g(j,x)>h_0\}}\right).
\]
To bound its expectation, expand this product. For each selected \((j,x)\) with \(g(j,x)=s>h_0\), bound by specifying \(s\) rows having \(X_{ij}=x\), summing over the choices and using factor \(e^{\gamma s}\) (for the probability bound we do not need to require that there are no further matching entries). In a term with several selected \((j,x)\), if the specifications conflict they have probability zero, and otherwise they require distinct cells. Thus
\[
 \mathbb E_{U_b^{\otimes a}} e^{\gamma R_X}
 \le \left(1+\sum_{s>h_0}\binom{a}{s} (a^{1/16} e/b)^s\right)^{b^2}
 \le e
\]
for all sufficiently large \(a\). Indeed the sum inside is at most \(\sum_{a\ge s>h_0} (2 e^2 a^{1/16}/s)^s\) using \(a/b\le 2\); this goes to zero faster than any inverse power of \(a\). Applying \eqref{tra:fractional-density:eq13} proves \eqref{tra:fractional-density:eq18}.

Since \(\epsilon\log a\) tends to infinity, \eqref{tra:fractional-density:eq17}-\eqref{tra:fractional-density:eq18} in \eqref{tra:fractional-density:eq15} absorb the terms \(\mathbb E_\zeta[R_X+R_Y]\) from \eqref{tra:fractional-density:eq16}. We conclude
\[
 D(\zeta\|\nu)
 \ \ge\ n-O\left(n a^{-1/16}\log a+(a+b)\log d\right).
\]
Taking \(\zeta=\nu(\cdot\mid\mathcal E)\) when \(\nu(\mathcal E)>0\), this proves \eqref{tra:fractional-density:eq14}.
\end{proof}

\paragraph{Completion of the rough moment estimate.}
Stirling's formula gives
\[
 \log\frac{n!}{|H||V|}
 =\log n!-a\log(b!)-b\log(a!)
 = n + O((a+b)\log n).
\]
Hence \eqref{tra:fractional-density:eq10}, \eqref{tra:fractional-density:eq14}, and \(M_b(p_0),M_a(p_0)\le 2\) bound \(M_n(p_0)\) by the exponential of at most \(O(n^{31/32}\log n)\), using \(a,b=\Theta(\sqrt{n})\). This proves \eqref{tra:fractional-density:eq7} for all sufficiently large \(d\).

\subsection{An infinity estimate on harmonic tuples}

\begin{proposition}\label{tra:fractional-density:harmonic}\label{e:harmonic-smoothing}

\label{tra:fractional-density:harmonic-infinity-statement}
Let $n=2^d$ with $d\ge1$, let $T_n$ be one coordinate sweep, and put $Q_n=T_n^*T_n$. Take \(1\le k\le n/(2d)\) and consider the representation on functions on ordered injective \(k\)-tuples of positions, with \(g\in S_n\) acting by
\[
 (\pi_k(g)F)(v)=F(g^{-1}v)
\]
where the notation on \(v\) acts on all its positions. Write \(\pi_k\) also for the corresponding action of group algebra elements. Call a function harmonic here if the sum in any one coordinate with the others fixed is 0 (taking the sum over the remaining available positions for that coordinate). The space of these functions is invariant under the permutation actions. We claim an infinity norm bound on this space:
\begin{equation}
 \|\pi_k(Q_n) F\|_\infty
 \ \le\ \left(\frac{C_* d k}{n}\right)^{k/2}\|F\|_\infty
 \qquad(F\ \text{harmonic})
 \label{tra:fractional-density:eq19}
\end{equation}
with $C_*=2(4e)^2$.

\end{proposition}\begin{proof}\label{tra:fractional-density:harmonic-proof}

\label{tra:fractional-density:midpoint-occupation-and-signed-path-laws}
In fact \((\pi_k(Q_n)F)(v)\) is the expectation of \(F\) on the tuple after moving \(v\) through a scan and then an independent reverse scan (layers in reverse order). This follows from \(Q_n=T_n^*T_n\) and the representation action. Start the tuple at \(v\) and let its ordered midpoint tuple between the scans be \(w\). For any set \(D\) of \(k\) positions, we have
\begin{equation}
 \Pr(\{w_1,\ldots,w_k\}=D)\le (k/n)^k .
 \label{tra:fractional-density:eq20}
\end{equation}
To see this, during the first scan consider just the occupied set of the \(k\) labels. At any point its law, with occupancy marginal probabilities \(p_x\) on positions, satisfies the bound \(\Pr(D'\ \text{all occupied})\le\prod_{x\in D'}p_x\) for any subset \(D'\). This holds initially, and is preserved through any one independent fair pair swap: the marginals of the two positions are replaced by their average. For a subset containing exactly one of these positions, averaging the two previous bounds gives the new bound; for a subset containing both, the probability stays the same and the product bound can only increase; for subsets containing neither there is no change. We can apply this through the sequence of all individual swaps. At the midpoint each card separately is uniform on the positions, since in its path each bit has been randomized with a fresh fair choice. Therefore \(p_x=k/n\) for the occupancy marginals, giving \eqref{tra:fractional-density:eq20}.

Condition on \(w\). Index the \(k\) labels in the tuple by \([k]\). For \(S\subseteq[k]\), consider a law \(\mathcal P_{S,w}\) of \(k\) paths in the second (reverse) scan with the following dynamics:
\begin{itemize}
\item the labels in \(S\) follow a common reverse scan with its random switches;
\item every other label (omitted from \(S\)) just follows a path using its own independent choices to stay or swap in the scheduled bit at each stage.

\end{itemize}
We allow the path positions under these laws not to form injective tuples across all the labels; extend \(F\) by 0 on noninjective tuples. For \(S\ne[k]\), the expectation of \(F\) at the endpoint under \(\mathcal P_{S,w}\) is 0. Indeed any omitted label has a uniform final position independent of the others. If the other final positions are distinct we use harmonicity and the extension by 0, and if not, \(F\) is already 0. Thus applying the signed sum of path laws
\begin{equation}
 \sum_{S\subseteq[k]} (-1)^{k-|S|}\mathcal P_{S,w}
 \label{tra:fractional-density:eq21}
\end{equation}
to \(F\) at the endpoint gives the same expectation as the true second scan.

\label{tra:fractional-density:isolated-cancellation-and-rooted-forest-count}
For any tuple of paths in these laws consider its encounter graph on \([k]\), joining two labels if they enter the same switch in some stage (possibly at the same position). The signed sum \eqref{tra:fractional-density:eq21} vanishes on any path tuple whose encounter graph has an isolated vertex. In fact if a vertex is isolated, toggling it in or out of \(S\) doesn't change the path tuple probability. To realize fixed paths in the common switches just requires the corresponding values of the switch coins, with probability zero in case of conflict. The isolated path uses \(d\) switches disjoint from those of any other paths (indexing switches by stage as well as pair), so it imposes its values independently, with probability \(2^{-d}\), just as if using its own omitted-label choices. This proves the cancellation.

Consequently the absolute conditional endpoint expectation in the true second scan is at most \(\|F\|_\infty G(w)\), where
\[
 G(w)=\sum_{S\subseteq[k]}\mathcal P_{S,w}(\text{encounter graph has no isolated vertex}).
\]
The function \(G(w)\) is symmetric under reordering \(w\): all choices of \(S\) are summed and the dynamics and condition correspond under renaming labels. Therefore \eqref{tra:fractional-density:eq20} gives
\begin{equation}
 \mathbb E G(w)\le \frac{(k/n)^k}{k!}\sum_{w\ \text{injective ordered}} G(w).
 \label{tra:fractional-density:eq22}
\end{equation}

To bound the last sum, fix \(S\) and fix its entire set of common switch coins and the separate stay/swap choice coins for omitted labels, using the same coins simultaneously for any starting \(w\). For each individual label, the map from its own starting position to its position entering any given stage is then a permutation of the \(n\) positions. This is true for a label in \(S\) by the common switch layers; for any other label the fixed choices apply bit flips or not, also giving permutations. For these coins, count the starting \(w\)'s whose encounter graph has no isolated vertex. Each such graph contains a spanning forest with \(c\le k/2\) components each of size at least 2; take a root in each component. (If \(k=1\) there are no such graphs.) For given \(c\), the number of rooted forests involved is at most
\(\binom{k}{c} k^{k-c}\), by choosing roots and parents. For a fixed forest, the starting positions of its roots can be chosen in at most \(n^c\) ways. For every other vertex, once the parent's starting position has been chosen, there are at most \(2d\) choices, since at some stage the child's path must enter the switch that the parent enters. At a given stage there are at most two choices of the child's starting position by the permutation property. Hence the number of \(w\)'s satisfying all the encounters of this forest is at most \(n^c(2d)^{k-c}\), not even needing to enforce injectivity in this bound. Using this count and averaging over coins for each \(S\), \eqref{tra:fractional-density:eq22} implies
\[
 \begin{aligned}
 \mathbb E G(w)
 &\le \frac{(k/n)^k}{k!}\,2^k
  \sum_{1\le c\le k/2} \binom{k}{c} (2d k)^{k-c} n^c\\
 &\le (2e)^k \sum_{1\le c\le k/2} \binom{k}{c} (2d k/n)^{k-c}
 \ \le\ \left(\frac{C_* d k}{n}\right)^{k/2}.
 \end{aligned}
\]
In the last inequality we used \(2dk/n\le 1\) and can take \(C_*=2(4e)^2\). This establishes \eqref{tra:fractional-density:eq19}.

\end{proof}

\begin{corollary}[Singular values at a first-row level]\label{e:harmonic-singular-conversion}
Let $\lambda\vdash n$ have first row $n-k$, where $1\le k\le n/(2d)$.
Then
\[
 \|Q_n(\lambda)\|_{\mathrm{op}}
 \le \left(\frac{2(4e)^2dk}{n}\right)^{k/2},
 \qquad
 \|T_n(\lambda)\|_{\mathrm{op}}
 \le \left(\frac{2(4e)^2dk}{n}\right)^{k/4}.
\]
\end{corollary}
\begin{proof}
The ordered $k$-tuple module is induced from the trivial representation of
$S_{n-k}$. Branching and Frobenius reciprocity show that $V_\lambda$
occurs in it, by removing the $k$ boxes below the first row. It does not
occur in any $(k-1)$-tuple module: every diagram there has first row at
least $n-k+1$. The span of functions lifted by omitting one coordinate
therefore has no $\lambda$-isotypic component. Its orthogonal complement
is exactly the harmonic space, because orthogonality to every such lift
means that each coordinate sum vanishes.

The operator $Q_n(\lambda)$ is positive. Every eigenvector in a copy of
$V_\lambda$ inside the harmonic space satisfies the infinity-norm bound
of Proposition~\ref{e:harmonic-smoothing}; dividing by its nonzero
infinity norm bounds its eigenvalue. This proves the first inequality.
The second follows from $Q_n(\lambda)=T_n(\lambda)^*T_n(\lambda)$.
\end{proof}
\subsection{Closing the trace recursion}

\begin{proof}[Proof of Proposition~\ref{tra:fractional-density:step}]

\label{tra:fractional-density:sparse-irrep-transfer}
Put $s=\lfloor n^{1-\beta}\rfloor$. For large $d$,
$s<n/(2d)$, so Corollary~\ref{e:harmonic-singular-conversion}
applies to all first-row levels $1\le k\le s$. Its positive-eigenvalue
bound is
\begin{equation}
 (C_*dk/n)^{k/2}\le n^{-\beta k/4},
 \label{tra:fractional-density:eq23}
\end{equation}
since $k/n\le n^{-\beta}$. There are at most $2^k$ shapes of level
$k$, each of dimension $r_\lambda\le n^k$. Thus their complete
regular contribution at any $p\ge16/\beta$ is at most
\begin{equation}
 \sum_{1\le k\le s}2^kn^{2k-p\beta k/4}=o(1).
 \label{tra:fractional-density:eq24}
\end{equation}
This uses the bound at the smallest allowed $p$, so the size threshold
is uniform over the inherited exponent.

\label{tra:fractional-density:sign-annihilation-and-catalan-dimensions}
A matching layer annihilates every shape with more than $n/2$
rows, as in the proof of Theorem~\ref{tra:duality:moment}.
In particular all shapes with $n-\lambda'_1\le s$ vanish, since
$n-s>n/2$ for sufficiently large $n$.

Treat the one-row shape separately: it is the trivial representation and contributes 1 to \(M_n(p)\). All shapes not handled so far have
\[
 n-\lambda_1>s,\qquad n-\lambda'_1>s.
\]
For them we have
\begin{equation}
 r_\lambda\ge \exp(c s)
 \label{tra:fractional-density:eq25}
\end{equation}
for some absolute \(c>0\) when \(d\) is sufficiently large. For if \(\lambda_1\ge s\), we can find a subdiagram consisting of a first row of length \(s\) and \(s\) more boxes below it forming a tail partition (as there are at least \(s\) boxes available below, taking a subpartition of size \(s\), of width necessarily at most \(s\)). The number of standard tableaux of this subdiagram is at least the Catalan number \(\frac{1}{s+1}\binom{2s}{s}\). Indeed interleave filling the top row and filling the \(s\)-box tail in the order of any fixed standard tableau on the tail, keeping the number of filled top row boxes always at least the number of filled tail boxes. This ballot condition ensures any tail box being filled already has the required top-row box above in its column (the column index is at most the number reached in filling the tail). Each resulting tableau extends to the original shape by filling the rest after the subdiagram. This proves \eqref{tra:fractional-density:eq25} in the case considered; the case \(\lambda'_1\ge s\) follows as well by transposing. If both lengths are smaller than \(s\), then every hook length is at most \(2s\), so
\[
 r_\lambda\ge n!/(2s)^n,
\]
which also gives \eqref{tra:fractional-density:eq25} for large enough \(d\), since \(n/(2s)\) tends to infinity.

\label{tra:fractional-density:large-dimension-trace-attenuation}
Apply the one-step estimate of Lemma~\ref{lem:balanced-moment-closure}
with the positive-square convention $\nu=2$, remainder $\varepsilon=1/8$,
low-level baseline $P=16/\beta$, and
\[
 E_n=n^{1-\alpha},\qquad H_n=cs,\qquad
 \alpha_d=1+d^{-1/2},\qquad\delta_d=d^{-1/2}.
\]
Here \eqref{tra:fractional-density:eq7} is the rough child-to-parent
bound, \eqref{tra:fractional-density:eq24} makes the low-level
remainder at most $1/16$, and \eqref{tra:fractional-density:eq25}
provides the regular multiplicity of every remaining block. The
large-block remainder is bounded by
\[
 \exp\{(1+d^{-1/2})n^{1-\alpha}
               -c d^{-1/2}\lfloor n^{1-\beta}\rfloor\}=o(1),
\]
because $\beta<\alpha$. It is therefore at most $1/16$ above an
absolute threshold independent of $q$. Adding the constant contribution
proves \eqref{tra:fractional-density:eq6}.

\end{proof}

\begin{theorem}\label{tra:fractional-density:moment}

\label{tra:fractional-density:bounded-moment-conclusion}
There are numbers \(p(d)>0\), with \(\sup_d p(d)<\infty\), such that for every \(n=2^d\),
\begin{equation}
 M_n(p(d))\le 1+c_0,\qquad c_0=\frac18.
 \label{tra:fractional-density:eq2}
\end{equation}

\end{theorem}
\begin{proof}[Proof of Theorem~\ref{tra:fractional-density:moment}]

\label{tra:fractional-density:finite-base-and-bounded-order}
The application of Lemma~\ref{lem:balanced-moment-closure} just
verified has exponent multiplier
$(1+d^{-1/2})^2=1+O(d^{-1/2})$. The lemma supplies the common finite
base and bounds the product of these multipliers over rounded
halvings. It therefore gives the asserted bounded sequence $p(d)$.
Proposition~\ref{e:nonnormal-moment-conversion} applies with singular
exponent $2\sup_dp(d)$ and remainder $1/8$.

\end{proof}
